% 原书物理页=印刷页16--19；2.3 全节.
\section{自由群与群的表述}\label{sec:fg-2-3}
\begin{numpara}[自由群]\label{np:fg-2-3-1}
设 $X$ 是一个非空集合. 若 $k\in\mathbb{N}$，字母表 $X$ 上一个长度为 $k$ 的\newterm{fg-word}{字}{word}是一个映射 $\{1,\cdots,k\}\longrightarrow X$. 若 $k=0$，这样的映射唯一（因为此时源是空集），称之为空字，记作 $e$. 若 $k>0$，一个字是 $X$ 中元素的一个序列 $m=(x_1x_2\cdots x_k)$. 记 $M(X)$ 为字母表 $X$ 上全体字的集合，并按下述方式赋予它一个内部运算. 对每个 $m$ 规定 $em=me=m$（故 $e$ 是单位元）. 对长度分别为 $k,k'$ 的 $m,m'$，规定 $mm'$ 为由拼接得到的长度 $k+k'$ 的字；换言之，$(x_1x_2\cdots x_k)(x_1'x_2'\cdots x_{k'}'):=(x_1x_2\cdots x_k\,x_1'x_2'\cdots x_{k'}')$. 这个运算显然是结合的，从而使 $M(X)$ 成为一个结合的\newterm{fg-monoid}{幺半群}{monoid}. 注意 $X$ 自然地等同于 $M(X)$ 中长度为 $1$ 的字所构成的子集. 容易验证，这个幺半群具有如下泛性质：任何从 $X$ 到一个结合幺半群的映射都唯一地延拓为一个幺半群同态 $M(X)\longrightarrow M$. 我们说 $M(X)$ 是集合 $X$ 上的\newterm{fg-free-monoid}{自由幺半群}{free monoid}.
\end{numpara}
注意，$M(X)$ 中任何非单位的元素都没有逆元；$M(X)$ 因此远不是一个群. 但我们现在想要造出一个具有与上述类似的泛性质的群 $F(X)$. 为此，形式地引入一个集合 $X^{-1}$，它配备一个从 $X$ 到 $X^{-1}$ 的双射，记作 $x\mapsto x^{-1}$. 于是考察字母表 $X\sqcup X^{-1}$（不相交并）上的字所成的幺半群 $M(X\sqcup X^{-1})$. 在 $M(X\sqcup X^{-1})$ 上定义一个序关系：称 $m\le m'$，如果从 $m'$ 出发，逐步消去形如 $xx^{-1}$ 或 $x^{-1}x$（$x\in X$）的子字便能得出 $m$. 称一个字为\newterm{fg-reduced-word}{既约字}{reduced word}，如果它关于这个序关系是极小的，即它不含形如 $xx^{-1}$ 或 $x^{-1}x$（$x\in X$）的子字. 对长度作归纳可证，每个字在序 $\le$ 下都大于或等于唯一的既约字. 由此得到一个从 $M(X\sqcup X^{-1})$ 到既约字所成子集 $F(X)$ 的``约化''映射. 再给 $F(X)$ 赋予如下运算：先拼接两个字再约化，即 $(m,m')\mapsto\operatorname{red}(mm')$. 由这个序关系与约化映射的性质可知，该运算仍是结合的，且空字是其单位元. 此外，每个元素都有逆元：$x\in X$（相应地，$x^{-1}$）的逆元是 $x^{-1}$（相应地，$x$）；任意字 $y_1\ldots y_n$（其中 $y_i\in X\cup X^{-1}$ 对每个 $i$ 成立）的逆元是 $y_n^{-1}\ldots y_1^{-1}$. 这样，$F(X)$ 是一个群，显然由它的子集 $X$ 生成.
\begin{proposition*}
对任意群 $G$，映射
\[
\operatorname{Hom}_{\mathrm{gpes}}(F(X),G)\longrightarrow G^X=\operatorname{Hom}_{\mathrm{ens}}(X,G)
\]
把 $\varphi:F(X)\longrightarrow G$ 对应到 $\varphi|_X$，是一个双射.
\end{proposition*}
换言之，任何从 $X$ 到群 $G$ 的映射都唯一地延拓为从 $F(X)$ 到 $G$ 的一个同态. 群 $F(X)$ 称为 $X$ 上的\newterm{fg-free-group}{自由群}{free group}.
\begin{proof}
由于 $F(X)$ 由 $X$ 生成，如上的同态 $\varphi$ 由它在 $X$ 上的限制完全确定，故题述映射是单射. 为证满射，任取 $\alpha:X\longrightarrow G$，把它延拓到 $X\cup X^{-1}$ 上，规定 $\alpha(x^{-1}):=\alpha(x)^{-1}$. 所得的映射又可延拓为一个映射 $M(X\sqcup X^{-1})\longrightarrow G$，它把字 $y_1\cdots y_n$ 映到乘积 $\alpha(y_1)\cdots\alpha(y_n)$. 这个映射与拼接相容. 另外，一个字的像与它相应的既约字的像一致. 因此它诱导出所求的群同态 $F(X)\longrightarrow G$.
\end{proof}
\begin{numpara}[例]\label{np:fg-2-3-2}
若 $X=\varnothing$，便得到平凡群. 若 $X$ 有 $n$ 个元素，也把 $F(X)$ 记作 $F(n)$，称之为 $n$ 个元素上的自由群，或称``秩为 $n$''的自由群. 群 $F(1)$ 同构于 $\mathbb{Z}$. 群 $F(2)$ 已经``极其庞大''. 每个自由群的子群都是自由群，这是一个非平凡的定理. 此外，当 $n>1$ 时，$F(n)$ 含有秩任意大的自由群！
\end{numpara}
\begin{numpara}[群的表述]\label{np:fg-2-3-3}
设 $S$ 是一个集合，$R\subset F(S)$ 是 $S$ 上既约字所组成的一个集合. 记 $\langle S|R\rangle$ 为群 $F(S)$ 对由 $R$ 生成的正规子群（即 $F(S)$ 中包含 $R$ 的最小正规子群）所作的商群. 我们说 $\langle S|R\rangle$ 是一个表述群.
\end{numpara}
\begin{definition*}
设 $G$ 是一个群. $G$ 的一个\newterm{fg-presentation}{表述}{presentation}\footnote{也称为``用生成元与关系给出的表述''.}由如下数据组成：一个如上的偶对 $(S,R)$ 以及一个映射 $S\longrightarrow G$，使得由它导出的典范同态 $F(S)\longrightarrow G$（见前一命题）诱导出一个同构 $\langle S|R\rangle\xrightarrow{\ \sim\ }G$.
\end{definition*}
任何群都拥有一个表述. 具体地说，设 $S$ 是 $G$ 的一个生成子集，即满足 $\langle S\rangle=G$ 的子集. 由 $F(S)$ 的泛性质导出的同态 $\pi:F(S)\to G$（在 $S$ 上诱导恒等映射）便是满射，且 $\pi$ 诱导出同构 $F(S)/\operatorname{Ker}(\pi)\xrightarrow{\ \sim\ }G$. 于是只须取 $R$ 为 $\operatorname{Ker}(\pi)$ 的任意一个子集，使得 $F(X)$ 中包含 $R$ 的最小正规子群恰为 $\operatorname{Ker}(\pi)$.\footnote{译注：原文此处写作 $F(X)$，依上下文应为 $F(S)$.}
应当注意，这样的表述绝不唯一. 事实上，这个概念定义起来相对简单，内涵却十分微妙. 例如，下面这些（模糊的）问题至今仍是当前研究的课题：
\begin{itemize}
\item 给定 $G$，找出 $G$ 的一个``尽可能简单''的表述.
\item 给定 $(S,R)$，计算群 $\langle S|R\rangle$（估计它的阶、指数等）.
\end{itemize}
为说明第二个问题，读者可以验证：由两个生成元给出的表述 $\langle\{a,b\}|\{a^4,b^3,a^{-2}bab^{-1}\}\rangle$ 是 $3$ 阶循环群，而由三个生成元给出的表述 $\langle\{a,b,c\}|\{a^{-1}bab^{-2},\,b^{-1}cbc^{-2},\,c^{-1}aca^{-2}\}\rangle$ 则是……平凡群.
\begin{numpara}[例]\label{np:fg-2-3-4}
$(\{a\}\mid\{a^n\})$ 是 $C_n$ 的一个表述；$(\{a,b\}\mid\{a^n,b^2,abab\})$ 是 $D_n$ 的一个表述；$(\{a,b\}\mid\{a^4,a^2b^2,ab^{-1}ab\})$ 是 $H_8$ 的一个表述. 注意，在每种情形下都容易找到从所说的表述群到声称给出表述的那个群的满同态；要证明核为零则更难. 不过，对于上面这样的小群，一个简单的计数论证便足以奏效.
\end{numpara}
\begin{numpara}[例：对称群]\label{np:fg-2-3-5}
在 $S_n$ 中，对 $i=1,\cdots,n-1$，把置换 $(i,i+1)$ 记作 $s_i$. 我们知道这些 $s_i$ 生成 $S_n$. 它们满足下列关系：
\begin{itemize}
\item $s_i^2=e$（二次关系）；
\item 当 $|i-j|>1$ 时 $(s_is_j)^2=e$，并且 $(s_is_{i+1})^3=e$（辫关系）.
\end{itemize}
令 $S=\{s_i,\ i=1,\cdots,n-1\}$、$R_0=\{s_i^2,\ i=1,\cdots,n-1\}$、$R_1=\{(s_is_{i+1})^3,\ (s_is_j)^2\}$ 如上，可以证明 $(S|R_0\cup R_1)$ 是群 $S_n$ 的一个表述. 我们也可以从一个 $n$-循环和一个置换出发来表述 $S_n$，但这并不是最有意思的表述. 事实上，把上述关系稍加改动，就得到两类非常有意思且被广泛研究的群：辫群与 Coxeter 群.
\end{numpara}
\begin{numpara}[辫群]\label{np:fg-2-3-6}
沿用上面的记号，把 $B_n$ 定义为表述群 $\langle S|R_1\rangle$. 这是一个无限群，称为``\newterm{fg-braid-group}{辫群}{braid group}''，在理论物理中十分重要. 它的元素可以用``辫子''来表示，也就是用 $n$ 根线绳把两组相对放置的、各由 $n$ 个点排成的点列一一连接起来的方式（物理上的点与真实的线绳；参见维基百科条目``braid groups''）. 它还有一个严格的拓扑学解释：它是复平面上 $n$ 个互不相同的点的构型空间的``基本群''.
\end{numpara}
\begin{numpara}[Coxeter 群]\label{np:fg-2-3-7}
一个 \newterm{fg-coxeter-group}{Coxeter 群}{Coxeter group}是形如 $\langle S,R\rangle$ 的表述群，其中关系集 $R$ 具有形式 $R=\{(st)^{m_{st}},\ s,t\in S\}$，这里 $(m_{st})_{s,t\in S}$ 是一些整数（可以是无穷），满足：对一切 $s$ 有 $m_{ss}=1$，而对一切 $s\neq t$ 有 $m_{st}=m_{ts}\geqslant 2$. 于是特别地对每个 $s$ 都有 $s^2=e$；因此群 $W$ 由一些阶为 $2$ 的元素生成. 条件 $m_{st}=\infty$ 只是意味着不在 $s$ 与 $t$ 之间施加任何关系；条件 $m_{st}=2$ 则等价于要求 $s$ 与 $t$ 交换.

这些表述的特殊之处在于``无冗余''：事实上，可以证明映射 $S\longrightarrow\langle S,R\rangle$ 是单射，并且元素 $st$ 在 $\langle S,R\rangle$ 中的阶是 $m_{st}$（这并非显而易见，见上文诸例）.

对称群是 Coxeter 群的特例（取 $S=\{1,\cdots,n-1\}$，$m_{i,i+1}=m_{i+1,i}=3$，而当 $|j-i|>1$ 时 $m_{ij}=2$）. 二面体群也是：取 $S=\{s,t\}$ 且 $m_{st}=m_{ts}=m$ 便得到 $D_m$. 一般地说，Coxeter 群未必有限，即使 $S$ 与诸 $m_{st}$ 都是有限的.

Coxeter 本人确定了哪些 Coxeter 群是有限的. 为此，他对数据 $(S|R)$ 配备一个图：其顶点集为 $S$，两个顶点 $s,t$ 之间连有 $m_{st}-2$ 条边. 他进而列举了所有使相应的群为有限的有限图. 只要 $|S|\geqslant 3$，以 $S$ 为顶点集的这类图就只有有限多个.

前面已经看到，二面体群是欧几里得平面上正多边形的对称群. 更一般地，$|S|=n$ 的有限 Coxeter 群是 $n$ 维欧几里得空间中某些（半）正``多胞形''的对称群. 例如，当 $|S|=3$ 时，这些多胞形就是``Platonic 立体''（正四面体、立方体、正八面体、正十二面体、正二十面体）. 第一个的等距群同构于 $S_4$，其后两个的等距群同构于半直积 $\{\pm1\}^4\rtimes S_4$，而最后两个的等距群则是 $A_5$.
\end{numpara}
\begin{numpara}[反射群]\label{np:fg-2-3-8}
设 $K$ 是一个域. 称 $\mathrm{GL}_n(K)$ 的一个自同构 $P$ 为\newterm{fg-reflection}{反射}{reflection}，如果它可对角化、具有有限阶，并且 $1$ 是重数为 $n-1$ 的特征值（此时唯一另外的特征值是一个单位根）. 超平面 $\operatorname{Ker}(P-I_n)$ 是该反射的轴，而直线 $\operatorname{Im}(P-I_n)$ 是它的方向. 称子群 $G\subset\mathrm{GL}_n(K)$ 为一个\newterm{fg-reflection-group}{反射群}{reflection group}，如果它由反射生成.

当 $K=\mathbb{R}$ 时，一个反射就是一个超平面对称，因而阶为 $2$. 这时可以证明，$\mathrm{GL}_n(\mathbb{R})$ 中任何由反射组成的子群都拥有一个 Coxeter 表述. 反过来，任何 Coxeter 群 $\langle S,R\rangle$ 都可以看作 $\mathrm{GL}_{|S|}(\mathbb{R})$ 中由反射组成的子群. 如果只限于有限群，那么在这两个断言中甚至可以把 $\mathrm{GL}_n(\mathbb{R})$ 换成 $\mathrm{O}_n(\mathbb{R})$（正交群，见下文）. 换言之，可以要求这些反射关于某个欧几里得内积是``正交的''（即轴与方向正交）.
\end{numpara}
